% Generated from homeworks.typ. Typst is authoritative; regeneration overwrites this file.
\documentclass{qlnotes-export}

\title{Math 525: Probability Homeworks}
\subtitle{Typst-first worked solutions}
\author{Qiulin Fan}
\date{2026}

\begin{document}
\mainmatter
\chapter{Homework 4}\label{homework-4}

\section{Problem 1}\label{problem-1-4}

Let \(X\) be a random variable with values in \(\lbrack 0, + \infty\rbrack\) such that \({\mathbb{E}}\lbrack X\rbrack = 0\). Explain why \(X < \infty\) almost surely and show that \(X = 0\) almost surely.

\begin{proof}

Since \(X \geq 0\) and \({\mathbb{E}}\lbrack X\rbrack = 0 < \infty\), by Markov's inequality, for any \(t > 0\),

\[{\mathbb{P}}(X \geq t) \leq \frac{{\mathbb{E}}\lbrack X\rbrack}{t} = 0\]

Hence \({\mathbb{P}}(X \geq t) = 0\) for all \(t > 0\). In particular,

\[{\mathbb{P}}(X = + \infty) = \lim\limits_{t\rightarrow + \infty}{\mathbb{P}}(X \geq t) = \lim\limits_{t\rightarrow + \infty}0 = 0\]

that is, \(X < \infty\) almost surely.

Also notice that

\[\left\{ {X > 0} \right\} = \bigcup\limits_{n = 1}^{\infty}\left\{ {X \geq \frac{1}{n}} \right\}\]

and by countable subadditivity,

\[{\mathbb{P}}(X > 0) \leq \sum\limits_{n = 1}^{\infty}{\mathbb{P}}\left( {X \geq \frac{1}{n}} \right) = 0\]

Thus \({\mathbb{P}}(X > 0) = 0\). And since \(X\) takes values in \(\lbrack 0, + \infty\rbrack\), we have

\[{\mathbb{P}}(\left\{ {X = 0} \right\}) = {\mathbb{P}}(\left\{ {X \leq 0} \right\}) = 1 - {\mathbb{P}}(\left\{ {X > 0} \right\}) = 1\]

, and

This finishes the proof that \(X < \infty\) a.s. and \(X = 0\) a.s.

\end{proof}

\section{Problem 2}\label{problem-2-4}

Let \(X\) be a random variable with \({\mathbb{E}}\lbrack X\rbrack = 3\) and \({\mathbb{E}}\lbrack X^{2}\rbrack = 13\). Show that:

\[{\mathbb{P}}( - 2 \leq X \leq 8) \geq \frac{21}{25}\]

\begin{qlsolution}{Solution}

Compute the variance of \(X\):

\[Var(X) = {\mathbb{E}}\lbrack X^{2}\rbrack - ({\mathbb{E}}\lbrack X\rbrack)^{2} = 13 - 3^{2} = 4\]

Observe that

\[\left. {\mathbb{P}}( - 2 \leq X \leq 8) = {\mathbb{P}}( \middle| X - 3 \middle| \leq 5) \right.\]

By Chebyshev's inequality,

\[\left. {\mathbb{P}}( \middle| X - 3 \middle| \geq 5) \leq \frac{Var(X)}{5^{2}} = \frac{4}{25} \right.\]

Therefore,

\[\left. {\mathbb{P}}( \middle| X - 3 \middle| \leq 5) = 1 - {\mathbb{P}}( \middle| X - 3 \middle| \geq 5) \geq 1 - \frac{4}{25} = \frac{21}{25} \right.\]

Thus,

\[{\mathbb{P}}( - 2 \leq X \leq 8) \geq \frac{21}{25}\]

\end{qlsolution}

\section{Problem 3}\label{problem-3-4}

Let \(X,Y\) be two random variables such that \({\mathbb{P}}(Y = 1) = 1/5,{\mathbb{P}}(Y = 2) = 3/5\) and \({\mathbb{P}}(Y = 3) = 1/5\). In addition

\[\left. X \middle| \left\{ {Y = 1} \right\} \sim \text{Exp}(2),X \middle| \left\{ {Y = 2} \right\} \sim \text{Exp}(3)\text{and}\ X \mid \left\{ {Y = 3} \right\} = 7. \right.\]

Compute the moment generating function of \(X\).

\begin{qlsolution}{Solution}

Compute each conditional moment generating function:

For \(X \mid \left\{ {Y = 1} \right\} \sim \text{Exp}(2)\)

\[{\mathbb{E}}\lbrack e^{tX} \mid Y = 1\rbrack = \frac{2}{2 - t},\quad t < 2\]

For \(X \mid \left\{ {Y = 2} \right\} \sim \text{Exp}(3)\)

\[{\mathbb{E}}\lbrack e^{tX} \mid Y = 2\rbrack = \frac{3}{3 - t},\quad t < 3\]

For \(X \mid \left\{ {Y = 3} \right\} = 7\),

\[{\mathbb{E}}\lbrack e^{tX} \mid Y = 3\rbrack = e^{7t}\]

Then we use the law of total expectation conditioning on \(Y\): for any \(t\) such that the expectations below are finite, we have

\[\begin{matrix}
{M_{X}(t) = {\mathbb{E}}\lbrack e^{tX}\rbrack} & {= \sum\limits_{y = 1}^{3}{\mathbb{E}}\lbrack e^{tX} \mid Y = y\rbrack{\mathbb{P}}(Y = y)} \\
 & {= \frac{1}{5} \cdot \frac{2}{2 - t} + \frac{3}{5} \cdot \frac{3}{3 - t} + \frac{1}{5}e^{7t}} \\
 & {= \frac{2}{5(2 - t)} + \frac{9}{5(3 - t)} + \frac{1}{5}e^{7t},\quad t < 2}
\end{matrix}\]

So the moment generating function of \(X\) is

\[M_{X}(t) = \frac{2}{5(2 - t)} + \frac{9}{5(3 - t)} + \frac{1}{5}e^{7t},\quad t < 2\]

\end{qlsolution}

\section{Problem 4}\label{problem-4-4}

For any \(n \in {\mathbb{N}}\) with \(n \geq 1\) we set \(a_{n} = 1/\left( {2n^{2}} \right)\). Consider the sequence of random variables \(\left( X_{n} \right)_{n \in {\mathbb{N}}}\) with

\[X_{n} = \left\{ \begin{matrix}
{0,} & {\text{with probability}\ a_{n}} \\
{1,} & {\text{with probability}\ 1 - 2a_{n}} \\
{n^{2},} & {\text{with probability}\ a_{n}}
\end{matrix} \right.\]

Check if \(\left( X_{n} \right)\) converges in probability and if \(\left( X_{n} \right)\) converges in \(L^{1}\).

\begin{qlsolution}{Solution}

We first show that \(X_{n}\rightarrow 1\) in probability.

Let \(\varepsilon > 0\).\\
For \(n\) large enough s.t. \(n^{2} - 1 > \varepsilon\), \(\left. |X_{n} - 1 \middle| > \varepsilon \right.\) iff \(X_{n} = 0\) or \(X_{n} = n^{2}\). Therefore,

\[\left. {\mathbb{P}}( \middle| X_{n} - 1 \middle| > \varepsilon) = {\mathbb{P}}(X_{n} = 0) + {\mathbb{P}}(X_{n} = n^{2}) = a_{n} + a_{n} = 2a_{n} = \frac{1}{n^{2}} \right.\]

Since \(\frac{1}{n^{2}}\rightarrow 0\), we have:

\[\left. \lim\limits_{n\rightarrow\infty}\ {\mathbb{P}}( \middle| X_{n} - 1 \middle| > \varepsilon) = 0 \right.\]

Since \(\varepsilon > 0\) is arbitrary, we conclude that

\[X_{n}\overset{\mathbb{P}}{\rightarrow}1\]

We then show that \(X_{n}\) does not converge to \(1\) in \(L^{1}\).

We compute

\[\begin{matrix}
\left. {\mathbb{E}}\lbrack \middle| X_{n} - 1 \middle| \rbrack \right. & \left. = \middle| 0 - 1 \middle| a_{n} + \middle| 1 - 1 \middle| (1 - 2a_{n}) + \middle| n^{2} - 1 \middle| a_{n} \right. \\
 & {= a_{n} + (n^{2} - 1)a_{n}} \\
 & {= n^{2}a_{n} = \frac{1}{2}\operatorname{\rightarrow\not{}}0}
\end{matrix}\]

Hence \$\$X\_n \textbackslash not\textbackslash xrightarrow\{L\^{}1\} 1\$\$

\end{qlsolution}

\section{Problem 5}\label{problem-5-4}

Let \(X\) and \(Y\) be independent random variables with densities

\[\begin{matrix}
{f_{X}(x) = \{ 2x,} & {0 \leq x \leq 1,} \\
{0,} & {\text{otherwise}\quad f_{Y}(y) = \left\{ \begin{matrix}
{1/2,} & {0 \leq y \leq 2} \\
{0,} & \text{otherwise}
\end{matrix} \right.}
\end{matrix}\]

Find the distribution function of the sum \(Z = X + Y\).

\begin{qlsolution}{Solution}

Since \(X\) and \(Y\) are independent, the density of \(Z = X + Y\) is given by convolution:

\[f_{Z}(z) = \int_{- \infty}^{\infty}f_{X}(x)f_{Y}(z - x)\, dx\]

where we know

\[f_{X}(x) = 2x\mathbf{1}_{\lbrack 0,1\rbrack}(x),\quad f_{Y}(y) = \frac{1}{2}\mathbf{1}_{\lbrack 0,2\rbrack}(y)\]

Thus

\[\begin{matrix}
{f_{Z}(z)} & {= \int_{- \infty}^{\infty}2x \cdot \frac{1}{2}\mathbf{1}_{\lbrack 0,1\rbrack}(x)\mathbf{1}_{\lbrack 0,2\rbrack}(z - x)\, dx} \\
 & {= \int_{- \infty}^{\infty}x\ \ \mathbf{1}_{\lbrack 0,1\rbrack}(x)\ \ \mathbf{1}_{\lbrack 0,2\rbrack}(z - x)\, dx} \\
 & {= \int_{\max(0,z - 2)}^{\min(1,z)}x\, dx}
\end{matrix}\]

Compute this piecewise: If \(0 \leq z \leq 1\), then the interval is \(\lbrack 0,z\rbrack\), so

\[f_{Z}(z) = \int_{0}^{z}x\, dx = \frac{z^{2}}{2}\]

If \(1 \leq z \leq 2\), then the interval is \(\lbrack 0,1\rbrack\), so

\[f_{Z}(z) = \int_{0}^{1}x\, dx = \frac{1}{2}\]

If \(2 \leq z \leq 3\), then the interval is \(\lbrack z - 2,1\rbrack\), so

\[f_{Z}(z) = \int_{z - 2}^{1}x\, dx = \frac{1 - (z - 2)^{2}}{2}\]

Outside \(\lbrack 0,3\rbrack\), clearly \(f_{Z}(z) = 0\). Therefore,

\[f_{Z}(z) = \left\{ \begin{matrix}
{0,} & {z < 0,} \\
{\frac{z^{2}}{2},} & {0 \leq z \leq 1,} \\
{\frac{1}{2},} & {1 \leq z \leq 2,} \\
{\frac{1 - (z - 2)^{2}}{2},} & {2 \leq z \leq 3,} \\
{0,} & {z > 3}
\end{matrix} \right.\]

Now we integrate to get the cdf. For \(z < 0\), \(F_{Z}(z) = 0\).

For \(0 \leq z \leq 1\),

\[F_{Z}(z) = \int_{0}^{z}\frac{t^{2}}{2}\, dt = \frac{z^{3}}{6}\]

For \(1 \leq z \leq 2\),

\[F_{Z}(z) = F_{Z}(1) + \int_{1}^{z}\frac{1}{2}\, dt = \frac{1}{6} + \frac{z - 1}{2} = \frac{z}{2} - \frac{1}{3}\]

For \(2 \leq z \leq 3\),

\[F_{Z}(z) = F_{Z}(2) + \int_{2}^{z}\frac{1 - (t - 2)^{2}}{2}\, dt = \frac{2}{3} + \frac{z - 2}{2} - \frac{(z - 2)^{3}}{6}\]

And for \(z \geq 3\), \(F_{Z}(z) = 1\).

Hence the cdf of \(Z = X + Y\) is

\[F_{Z}(z) = \left\{ \begin{matrix}
{0,} & {z < 0,} \\
{\frac{z^{3}}{6},} & {0 \leq z \leq 1,} \\
{\frac{z}{2} - \frac{1}{3},} & {1 \leq z \leq 2,} \\
{\frac{2}{3} + \frac{z - 2}{2} - \frac{(z - 2)^{3}}{6},} & {2 \leq z \leq 3,} \\
{1,} & {z \geq 3}
\end{matrix} \right.\]

\end{qlsolution}

\section{Problem 6}\label{problem-6-4}

Let \(a_{1},a_{2},\ldots,a_{n}\) and \(\lambda\) be positive constants and let \(\left\{ {X_{i}:1 \leq i \leq n} \right\}\) be independent random variables with

\[X_{i} \sim \Gamma\left( {a_{i},\lambda} \right),\quad i = 1,2,\ldots,n\]

(i.e., with the same second parameter \(\lambda\) ). Show that \(X_{1} + X_{2} + \cdots + X_{n} \sim \Gamma(a,\lambda)\), with \(a = \sum_{i = 1}^{n}a_{i}\).

\begin{proof}

Let

\[S_{n} := X_{1} + X_{2} + \cdots + X_{n},\quad a = \sum\limits_{i = 1}^{n}a_{i}\]

We need to show that \(S_{n} \sim \Gamma(a,\lambda)\).

Since \(X_{i} \sim \Gamma(a_{i},\lambda)\), its moment generating function is

\[M_{X_{i}}(t) = {\mathbb{E}}\lbrack e^{tX_{i}}\rbrack = \left( \frac{\lambda}{\lambda - t} \right)^{a_{i}},\quad t < \lambda\]

Since \(X_{1},\ldots,X_{n}\) are independent, the moment generating function of \(S_{n}\) is

\[M_{S_{n}}(t) = {\mathbb{E}}\lbrack e^{t(X_{1} + \cdots + X_{n})}\rbrack = \prod\limits_{i = 1}^{n}{\mathbb{E}}\lbrack e^{tX_{i}}\rbrack = \prod\limits_{i = 1}^{n}M_{X_{i}}(t)\]

Therefore,

\[M_{S_{n}}(t) = \prod\limits_{i = 1}^{n}\left( \frac{\lambda}{\lambda - t} \right)^{a_{i}} = \left( \frac{\lambda}{\lambda - t} \right)^{\sum_{i = 1}^{n}a_{i}} = \left( \frac{\lambda}{\lambda - t} \right)^{a},\quad t < \lambda\]

Note this is exactly the moment generating function of a \(\Gamma(a,\lambda)\) random variable. Hence,

\[S_{n} = X_{1} + \cdots + X_{n} \sim \Gamma(a,\lambda),\quad a = \sum\limits_{i = 1}^{n}a_{i}\]

\end{proof}
\end{document}
